%=====================================
%   20250802.tex
%   20250802 
%   toshi2019
%=====================================
% ------------------------
% documentclass
% ------------------------
\documentclass[leqno, 10pt, a4paper, dvipdfmx]{jsarticle}
\title{ゼミノート（10/08分）}
\author{Toshi2019}
\date{\today}

% ------------------------
% usepackage
% ------------------------
\usepackage{algorithm}
\usepackage{algorithmic}
\usepackage{amscd}
\usepackage{amsfonts}
\usepackage{amsmath}
\usepackage[psamsfonts]{amssymb}
\usepackage{amsthm}
\usepackage{ascmac}
\usepackage{bm}
\usepackage{color}
\usepackage{enumerate}
\usepackage{fancybox}
\usepackage[stable]{footmisc}
\usepackage{graphicx}
\usepackage{listings}
\usepackage{mathrsfs}
\usepackage{mathtools}
\usepackage{otf}
\usepackage{pifont}
\usepackage{proof}
\usepackage{subfigure}
\usepackage{tikz}
\usepackage{verbatim}
\usepackage[all]{xy}

\usetikzlibrary{cd}
\usetikzlibrary{arrows.meta}
\usetikzlibrary{calc}

% ================================
% パッケージを追加する場合のスペース 
\usepackage[dvipdfmx]{hyperref}
\usepackage{xcolor}
\definecolor{darkgreen}{rgb}{0,0.45,0} 
\definecolor{darkred}{rgb}{0.75,0,0}
\definecolor{darkblue}{rgb}{0,0,0.6} 
\hypersetup{
    colorlinks=true,
    citecolor=darkgreen,
    linkcolor=darkred,
    urlcolor=darkblue,
}
\usepackage{pxjahyper}

%\usepackage[mathscr]{euscript} % mathscr のフォントを変更

%\usepackage{mmasym}

%=================================


% --------------------------
% theoremstyle
% --------------------------
\theoremstyle{definition}

% --------------------------
% newtheoem
% --------------------------

% 日本語で定理, 命題, 証明などを番号付きで用いるためのコマンドです. 
% If you want to use theorem environment in Japanece, 
% you can use these code. 
% Attention!
% All theorem enivironment numbers depend on 
% only section numbers.
\newtheorem{Axiom}{公理}[section]
\newtheorem{Definition}[Axiom]{定義}
\newtheorem{Theorem}[Axiom]{定理}
\newtheorem{Proposition}[Axiom]{命題}
\newtheorem{Lemma}[Axiom]{補題}
\newtheorem{Corollary}[Axiom]{系}
\newtheorem{Example}[Axiom]{例}
\newtheorem{Claim}[Axiom]{主張}
\newtheorem{Property}[Axiom]{性質}
\newtheorem{Attention}[Axiom]{注意}
\newtheorem{Question}[Axiom]{問}
\newtheorem{Problem}[Axiom]{問題}
\newtheorem{Consideration}[Axiom]{考察}
\newtheorem{Alert}[Axiom]{警告}
\newtheorem{Fact}[Axiom]{事実}


% 日本語で定理, 命題, 証明などを番号なしで用いるためのコマンドです. 
% If you want to use theorem environment with no number in Japanese, You can use these code.
\newtheorem*{Axiom*}{公理}
\newtheorem*{Definition*}{定義}
\newtheorem*{Theorem*}{定理}
\newtheorem*{Proposition*}{命題}
\newtheorem*{Lemma*}{補題}
\newtheorem*{Example*}{例}
\newtheorem*{Corollary*}{系}
\newtheorem*{Claim*}{主張}
\newtheorem*{Property*}{性質}
\newtheorem*{Attention*}{注意}
\newtheorem*{Question*}{問}
\newtheorem*{Problem*}{問題}
\newtheorem*{Consideration*}{考察}
\newtheorem*{Alert*}{警告}
\newtheorem{Fact*}{事実}


% 英語で定理, 命題, 証明などを番号付きで用いるためのコマンドです. 
%　If you want to use theorem environment in English, You can use these code.
%all theorem enivironment number depend on only section number.
\newtheorem{Axiom+}{Axiom}[section]
\newtheorem{Definition+}[Axiom+]{Definition}
\newtheorem{Theorem+}[Axiom+]{Theorem}
\newtheorem{Proposition+}[Axiom+]{Proposition}
\newtheorem{Lemma+}[Axiom+]{Lemma}
\newtheorem{Example+}[Axiom+]{Example}
\newtheorem{Corollary+}[Axiom+]{Corollary}
\newtheorem{Claim+}[Axiom+]{Claim}
\newtheorem{Property+}[Axiom+]{Property}
\newtheorem{Attention+}[Axiom+]{Attention}
\newtheorem{Question+}[Axiom+]{Question}
\newtheorem{Problem+}[Axiom+]{Problem}
\newtheorem{Consideration+}[Axiom+]{Consideration}
\newtheorem{Alert+}{Alert}
\newtheorem{Fact+}[Axiom+]{Fact}
\newtheorem{Remark+}[Axiom+]{Remark}

% ----------------------------
% commmand
% ----------------------------
% 執筆に便利なコマンド集です. 
% コマンドを追加する場合は下のスペースへ. 

% 集合の記号 (黒板文字)
\newcommand{\NN}{\mathbb{N}}
\newcommand{\ZZ}{\mathbb{Z}}
\newcommand{\QQ}{\mathbb{Q}}
\newcommand{\RR}{\mathbb{R}}
\newcommand{\CC}{\mathbb{C}}
\newcommand{\PP}{\mathbb{P}}
\newcommand{\KK}{\mathbb{K}}


% 集合の記号 (太文字)
\newcommand{\nn}{\mathbf{N}}
\newcommand{\zz}{\mathbf{Z}}
\newcommand{\qq}{\mathbf{Q}}
\newcommand{\rr}{\mathbf{R}}
\newcommand{\cc}{\mathbf{C}}
\newcommand{\pp}{\mathbf{P}}
\newcommand{\kk}{\mathbf{k}}

% 特殊な写像の記号
\newcommand{\ev}{\mathop{\mathrm{ev}}\nolimits} % 値写像
\newcommand{\pr}{\mathop{\mathrm{pr}}\nolimits} % 射影

% スクリプト体にするコマンド
%   例えば　{\mcal C} のように用いる
\newcommand{\mcal}{\mathcal}

% 花文字にするコマンド 
%   例えば　{\h C} のように用いる
\newcommand{\h}{\mathscr}

% ヒルベルト空間などの記号
\newcommand{\F}{\mcal{F}}
\newcommand{\X}{\mcal{X}}
\newcommand{\Y}{\mcal{Y}}
\newcommand{\Hil}{\mcal{H}}
\newcommand{\RKHS}{\Hil_{k}}
\newcommand{\Loss}{\mcal{L}_{D}}
\newcommand{\MLsp}{(\X, \Y, D, \Hil, \Loss)}

% 偏微分作用素の記号
\newcommand{\p}{\partial}

% 角カッコの記号 (内積は下にマクロがあります)
\newcommand{\lan}{\langle}
\newcommand{\ran}{\rangle}



% 圏の記号など
\newcommand{\Set}{\mathsf{Set}}
\newcommand{\Vect}{\mathsf{Vect}}
\newcommand{\FDVect}{\mathsf{FDVect}}
%\newcommand{\FDVect}{{\bf FDVect}}
%\newcommand{\Ring}{{\bf Ring}}
\newcommand{\Ab}{{\bf Ab}}
\newcommand{\Mod}{\mathop{\mathrm{Mod}}\nolimits}
\newcommand{\Modf}{\mathop{\mathrm{Mod}^\mathrm{f}}\nolimits}
\newcommand{\CGA}{{\bf CGA}}
\newcommand{\GVect}{{\bf GVect}}
\newcommand{\Lie}{{\bf Lie}}
\newcommand{\dLie}{{\bf Liec}}



% 射の集合など
\newcommand{\Map}{\mathop{\mathrm{Map}}\nolimits} % 写像の集合
\newcommand{\Hom}{\mathop{\mathrm{Hom}}\nolimits} % 射集合
\newcommand{\End}{\mathop{\mathrm{End}}\nolimits} % 自己準同型の集合
\newcommand{\Aut}{\mathop{\mathrm{Aut}}\nolimits} % 自己同型の集合
\newcommand{\Mor}{\mathop{\mathrm{Mor}}\nolimits} % 射集合
\newcommand{\Ker}{\mathop{\mathrm{Ker}}\nolimits} % 核
\newcommand{\Img}{\mathop{\mathrm{Im}}\nolimits} % 像
\newcommand{\Cok}{\mathop{\mathrm{Coker}}\nolimits} % 余核
\newcommand{\Cim}{\mathop{\mathrm{Coim}}\nolimits} % 余像

% その他便利なコマンド
\newcommand{\dip}{\displaystyle} % 本文中で数式モード
\newcommand{\e}{\varepsilon} % イプシロン
\newcommand{\dl}{\delta} % デルタ
\newcommand{\pphi}{\varphi} % ファイ
\newcommand{\ti}{\tilde} % チルダ
\newcommand{\pal}{\parallel} % 平行
\newcommand{\op}{{\rm op}} % 双対を取る記号
\newcommand{\lcm}{\mathop{\mathrm{lcm}}\nolimits} % 最小公倍数の記号
\newcommand{\Probsp}{(\Omega, \F, \P)} 
\newcommand{\argmax}{\mathop{\rm arg~max}\limits}
\newcommand{\argmin}{\mathop{\rm arg~min}\limits}





% ================================
% コマンドを追加する場合のスペース 
%\newcommand{\OO}{\mcal{O}}



\makeatletter
\renewenvironment{proof}[1][\proofname]{\par
  \pushQED{\qed}%
  \normalfont \topsep6\p@\@plus6\p@\relax
  \trivlist
  \item[\hskip\labelsep
%        \itshape
         \bfseries
%    #1\@addpunct{.}]\ignorespaces
    {#1}]\ignorespaces
}{%
  \popQED\endtrivlist\@endpefalse
}
\makeatother

\renewcommand{\proofname}{証明.}



%\renewcommand\proofname{\bf 証明} % 証明
\numberwithin{equation}{section}
\newcommand{\cTop}{\textsf{Top}}
%\newcommand{\cOpen}{\textsf{Open}}
\newcommand{\Op}{\mathop{\textsf{Op}}\nolimits}
\newcommand{\Ob}{\mathop{\textrm{Ob}}\nolimits}
\newcommand{\id}{\mathop{\mathrm{id}}\nolimits}
\newcommand{\pt}{\mathop{\mathrm{pt}}\nolimits}
\newcommand{\res}{\mathop{\rho}\nolimits}
\newcommand{\A}{\mcal{A}}
\newcommand{\B}{\mcal{B}}
\newcommand{\C}{\mcal{C}}
\newcommand{\D}{\mcal{D}}
\newcommand{\E}{\mcal{E}}
\newcommand{\G}{\mcal{G}}
%\newcommand{\H}{\mcal{H}}
\newcommand{\I}{\mcal{I}}
\newcommand{\J}{\mcal{J}}
\newcommand{\OO}{\mcal{O}}
\newcommand{\Ring}{\mathop{\textsf{Ring}}\nolimits}
\newcommand{\cAb}{\mathop{\textsf{Ab}}\nolimits}
%\newcommand{\Ker}{\mathop{\mathrm{Ker}}\nolimits}
\newcommand{\im}{\mathop{\mathrm{Im}}\nolimits}
\newcommand{\Coker}{\mathop{\mathrm{Coker}}\nolimits}
\newcommand{\Coim}{\mathop{\mathrm{Coim}}\nolimits}
\newcommand{\rank}{\mathop{\mathrm{rank}}\nolimits}
\newcommand{\Ht}{\mathop{\mathrm{Ht}}\nolimits}
\newcommand{\supp}{\mathop{\mathrm{supp}}\nolimits}
\newcommand{\colim}{\mathop{\mathrm{colim}}}
\newcommand{\Tor}{\mathop{\mathrm{Tor}}\nolimits}

\newcommand{\cat}{\mathcal{C}}

%筆記体
\newcommand{\cA}{\mcal{A}}
\newcommand{\cB}{\mcal{B}}
\newcommand{\cC}{\mcal{C}}
\newcommand{\cD}{\mcal{D}}
\newcommand{\cE}{\mcal{E}}
\newcommand{\cF}{\mcal{F}}
\newcommand{\cG}{\mcal{G}}
\newcommand{\cH}{\mcal{H}}
\newcommand{\cI}{\mcal{I}}
\newcommand{\cJ}{\mcal{J}}
\newcommand{\cK}{\mcal{K}}
\newcommand{\cL}{\mcal{L}}
\newcommand{\cM}{\mcal{M}}
\newcommand{\cN}{\mcal{N}}
\newcommand{\cO}{\mcal{O}}
\newcommand{\cP}{\mcal{P}}
\newcommand{\cQ}{\mcal{Q}}
\newcommand{\cR}{\mcal{R}}
\newcommand{\cS}{\mcal{S}}
\newcommand{\cT}{\mcal{T}}
\newcommand{\cU}{\mcal{U}}
\newcommand{\cV}{\mcal{V}}
\newcommand{\cW}{\mcal{W}}
\newcommand{\cX}{\mcal{X}}
\newcommand{\cY}{\mcal{Y}}
\newcommand{\cZ}{\mcal{Z}}


\newcommand{\scA}{\mathscr{A}}
\newcommand{\scB}{\mathscr{B}}
\newcommand{\scC}{\mathscr{C}}
\newcommand{\scD}{\mathscr{D}}
\newcommand{\scE}{\mathscr{E}}
\newcommand{\scF}{\mathscr{F}}
\newcommand{\scN}{\mathscr{N}}
\newcommand{\scO}{\mathscr{O}}
\newcommand{\scV}{\mathscr{V}}
\newcommand{\scU}{\mathscr{U}}


\newcommand{\ibA}{\mathop{\text{\textit{\textbf{A}}}}}
\newcommand{\ibB}{\mathop{\text{\textit{\textbf{B}}}}}
\newcommand{\ibC}{\mathop{\text{\textit{\textbf{C}}}}}
\newcommand{\ibD}{\mathop{\text{\textit{\textbf{D}}}}}
\newcommand{\ibE}{\mathop{\text{\textit{\textbf{E}}}}}
\newcommand{\ibF}{\mathop{\text{\textit{\textbf{F}}}}}
\newcommand{\ibG}{\mathop{\text{\textit{\textbf{G}}}}}
\newcommand{\ibH}{\mathop{\text{\textit{\textbf{H}}}}}
\newcommand{\ibI}{\mathop{\text{\textit{\textbf{I}}}}}
\newcommand{\ibJ}{\mathop{\text{\textit{\textbf{J}}}}}
\newcommand{\ibK}{\mathop{\text{\textit{\textbf{K}}}}}
\newcommand{\ibL}{\mathop{\text{\textit{\textbf{L}}}}}
\newcommand{\ibM}{\mathop{\text{\textit{\textbf{M}}}}}
\newcommand{\ibN}{\mathop{\text{\textit{\textbf{N}}}}}
\newcommand{\ibO}{\mathop{\text{\textit{\textbf{O}}}}}
\newcommand{\ibP}{\mathop{\text{\textit{\textbf{P}}}}}
\newcommand{\ibQ}{\mathop{\text{\textit{\textbf{Q}}}}}
\newcommand{\ibR}{\mathop{\text{\textit{\textbf{R}}}}}
\newcommand{\ibS}{\mathop{\text{\textit{\textbf{S}}}}}
\newcommand{\ibT}{\mathop{\text{\textit{\textbf{T}}}}}
\newcommand{\ibU}{\mathop{\text{\textit{\textbf{U}}}}}
\newcommand{\ibV}{\mathop{\text{\textit{\textbf{V}}}}}
\newcommand{\ibW}{\mathop{\text{\textit{\textbf{W}}}}}
\newcommand{\ibX}{\mathop{\text{\textit{\textbf{X}}}}}
\newcommand{\ibY}{\mathop{\text{\textit{\textbf{Y}}}}}
\newcommand{\ibZ}{\mathop{\text{\textit{\textbf{Z}}}}}

\newcommand{\ibx}{\mathop{\text{\textit{\textbf{x}}}}}

%\newcommand{\Comp}{\mathop{\mathrm{C}}\nolimits}
%\newcommand{\Komp}{\mathop{\mathrm{K}}\nolimits}
%\newcommand{\Domp}{\mathop{\mathsf{D}}\nolimits}%複体のホモトピー圏
%\newcommand{\Comp}{\mathrm{C}}
%\newcommand{\Komp}{\mathrm{K}}
%\newcommand{\Domp}{\mathsf{D}}%複体のホモトピー圏

\newcommand{\Comp}{\mathop{\mathrm{C}}\nolimits}
\newcommand{\Komp}{\mathop{\mathsf{K}}\nolimits}
\newcommand{\Domp}{\mathord{\mathsf{D}}}%\nolimits}
\newcommand{\Kompl}{\mathop{\mathsf{K}^\mathrm{+}}\nolimits}
\newcommand{\Kompu}{\mathop{\mathsf{K}^\mathrm{-}}\nolimits}
\newcommand{\Kompb}{\mathop{\mathsf{K}^\mathrm{b}}\nolimits}
\newcommand{\Dompl}{\mathop{\mathsf{D}^\mathrm{+}}\nolimits}
\newcommand{\Dompu}{\mathop{\mathsf{D}^\mathrm{-}}\nolimits}
\newcommand{\Dompb}{\mathop{\mathsf{D}^\mathrm{b}}\nolimits}
\newcommand{\Dbconic}{\mathop{\mathsf{D}^\mathrm{b}_{\rrp}}\nolimits}
\newcommand{\Dompbf}{\mathop{\mathsf{D}_\mathrm{f}^\mathrm{b}}\nolimits}
\newcommand{\Domplb}{\mathop{\mathsf{D}^\mathrm{lb}}\nolimits}




\newcommand{\CCat}{\Comp(\cat)}
\newcommand{\KCat}{\Komp(\cat)}
\newcommand{\DCat}{\Domp(\cat)}%圏Cの複体のホモトピー圏
%\newcommand{\HOM}{\mathop{\mathscr{H}\hspace{-2pt}om}\nolimits}%内部Hom
\newcommand{\HOM}{\mathop{\mathcal{H}\hspace{-0pt}om}\nolimits}%内部Hom
\newcommand{\RHOM}{\mathop{\mathrm{R}\hspace{-1.5pt}\HOM}\nolimits}

\newcommand{\muS}{\mathop{\mathrm{SS}}\nolimits}
\newcommand{\RG}{\mathop{\mathrm{R}\hspace{-0pt}\Gamma}\nolimits}
\newcommand{\RHom}{\mathop{\mathrm{R}\hspace{-1.5pt}\Hom}\nolimits}
\newcommand{\Rder}{\mathrm{R}}

\newcommand{\simar}{\mathrel{\overset{\sim}{\rightarrow}}}%同型右矢印
\newcommand{\simarr}{\mathrel{\overset{\sim}{\longrightarrow}}}%同型右矢印
\newcommand{\simra}{\mathrel{\overset{\sim}{\leftarrow}}}%同型左矢印
\newcommand{\simrra}{\mathrel{\overset{\sim}{\longleftarrow}}}%同型左矢印

\newcommand{\hocolim}{{\mathrm{hocolim}}}
\newcommand{\indlim}[1][]{\mathop{\varinjlim}\limits_{#1}}
\newcommand{\sindlim}[1][]{\smash{\mathop{\varinjlim}\limits_{#1}}\,}
\newcommand{\Pro}{\mathrm{Pro}}
\newcommand{\Ind}{\mathrm{Ind}}
\newcommand{\prolim}[1][]{\mathop{\varprojlim}\limits_{#1}}
\newcommand{\sprolim}[1][]{\smash{\mathop{\varprojlim}\limits_{#1}}\,}

\newcommand{\Sh}{\mathrm{Sh}}
\newcommand{\PSh}{\mathrm{PSh}}

\newcommand{\rmD}{\mathsf{D}}
\newcommand{\vD}{\mathbf{D}}

\newcommand{\Lloc}[1][]{\mathord{\mathcal{L}^1_{\mathrm{loc},{#1}}}}
\newcommand{\ori}{\mathord{\mathrm{or}}}
\newcommand{\Db}{\mathord{\mathscr{D}b}}

\newcommand{\codim}{\mathop{\mathrm{codim}}\nolimits}



\newcommand{\gld}{\mathop{\mathrm{gld}}\nolimits}
\newcommand{\wgld}{\mathop{\mathrm{wgld}}\nolimits}


\newcommand{\tens}[1][]{\mathbin{\otimes_{\raise1.5ex\hbox to-.1em{}{#1}}}}
\newcommand{\etens}{\mathbin{\boxtimes}}
\newcommand{\ltens}[1][]{\mathbin{\overset{\mathrm{L}}\tens}_{#1}}
\newcommand{\mtens}[1][]{\mathbin{\overset{\mathrm{\mu}}\tens}_{#1}}
\newcommand{\lltens}[1][]{\mathbin{{\mathop{\tens}\limits^{\mathrm{L}}_{#1}}}}
\newcommand{\lletens}[1][]{\mathbin{{\mathop{\boxtimes}\limits^{\mathrm{L}}_{#1}}}}
\newcommand{\letens}{\overset{\mathrm{L}}{\etens}}
\newcommand{\detens}{\underline{\etens}}
\newcommand{\ldetens}{\overset{\mathrm{L}}{\underline{\etens}}}
\newcommand{\dtens}[1][]{{\overset{\mathrm{L}}{\underline{\otimes}}}_{#1}}

\newcommand{\blk}{\mathord{\ \cdot\ }}
\newcommand{\mres}[2][]{{\left.{#1}\right\rvert}_{#2}}


%\newcommand{\hocolim}{{\mathrm{hocolim}}}
%\newcommand{\indlim}[1][]{\mathop{\varinjlim}\limits_{#1}}
%\newcommand{\sindlim}[1][]{\smash{\mathop{\varinjlim}\limits_{#1}}\,}
%\newcommand{\Pro}{\mathrm{Pro}}
%\newcommand{\Ind}{\mathrm{Ind}}
%\newcommand{\prolim}[1][]{\mathop{\varprojlim}\limits_{#1}}
%\newcommand{\sprolim}[1][]{\smash{\mathop{\varprojlim}\limits_{#1}}\,}
\newcommand{\proolim}[1][]{\mathop{\text{\rm``{$\varprojlim$}''}}\limits_{#1}}
\newcommand{\sproolim}[1][]{\smash{\mathop{\rm``{\varprojlim}''}\limits_{#1}}}
\newcommand{\inddlim}[1][]{\mathop{\text{\rm``{$\varinjlim$}''}}\limits_{#1}}
\newcommand{\sinddlim}[1][]{\smash{\mathop{\text{\rm``{$\varinjlim$}''}}\limits_{#1}}\,}
\newcommand{\ooplus}{\mathop{\text{\rm``{$\oplus$}''}}\limits}
\newcommand{\bbigsqcup}{\mathop{``\bigsqcup''}\limits}
\newcommand{\bsqcup}{\mathop{``\sqcup''}\limits}
\newcommand{\dsum}[1][]{\mathbin{\oplus_{#1}}}

\newcommand{\Fct}{\mathop{\mathsf{Fct}}\nolimits}
%Yoneda embedding
\newcommand{\yoneda}[1][]{\mathrm{h}_{#1}}
\newcommand{\coyoneda}[1][]{\mathrm{k}_{#1}}
%presheaf category
\newcommand{\precat}[1][]{{#1}^{\wedge}}
\newcommand{\coprecat}[1][]{{#1}^{\vee}}

\newcommand{\pb}[1][]{\mathbin{\mathop{\times}\limits_{#1}}}
\newcommand{\dpb}[1][]{\mathbin{\mathop{\times}\limits_{#1}}}
\newcommand{\tpb}[1][]{\mathbin{\mathop{\times}\nolimits_{#1}}}

\newcommand{\eu}{\mathbf{e}} % the Euler vector field
\newcommand{\rrp}{\rr^{+}} % positive real numbers




%================================================
% 自前の定理環境
%   https://mathlandscape.com/latex-amsthm/
% を参考にした
\newtheoremstyle{mystyle}%   % スタイル名
    {5pt}%                   % 上部スペース
    {5pt}%                   % 下部スペース
    {}%              % 本文フォント
    {}%                  % 1行目のインデント量
    {\bfseries}%                      % 見出しフォント
    {.}%                     % 見出し後の句読点
    {12pt}%                     % 見出し後のスペース
    {\thmname{#1}\thmnumber{ #2}\thmnote{{\hspace{2pt}\normalfont (#3)}}}% % 見出しの書式

\theoremstyle{mystyle}
\newtheorem{AXM}{公理}[section]
\newtheorem{DFN}[AXM]{定義}
\newtheorem{THM}[AXM]{定理}
\newtheorem*{THM*}{定理}
\newtheorem{PRP}[AXM]{命題}
\newtheorem{LMM}[AXM]{補題}
\newtheorem{CNJ}[AXM]{予想}
\newtheorem{CRL}[AXM]{系}
\newtheorem{EG}[AXM]{例}
\newtheorem*{EG*}{例}
\newtheorem{RMK}[AXM]{注意}
\newtheorem{CNV}[AXM]{約束}
\newtheorem{CMT}[AXM]{コメント}
\newtheorem*{CMT*}{コメント}
\newtheorem{NTN}[AXM]{記号}

% 定理環境ここまで
%====================================================

\usepackage{framed}
\definecolor{lightgray}{rgb}{0.75,0.75,0.75}
\renewenvironment{leftbar}{%
  \def\FrameCommand{\textcolor{lightgray}{\vrule width 4pt} \hspace{10pt}}% 
  \MakeFramed {\advance\hsize-\width \FrameRestore}}%
{\endMakeFramed}
\newenvironment{redleftbar}{%
  \def\FrameCommand{\textcolor{lightgray}{\vrule width 1pt} \hspace{10pt}}% 
  \MakeFramed {\advance\hsize-\width \FrameRestore}}%
 {\endMakeFramed}



\renewcommand{\Re}{\mathop{\mathrm{Re}}\nolimits}

% =================================





% ---------------------------
% new definition macro
% ---------------------------
% 便利なマクロ集です

% 内積のマクロ
%   例えば \inner<\pphi | \psi> のように用いる
\def\inner<#1>{\langle #1 \rangle}

% ================================
% マクロを追加する場合のスペース 

%=================================



% 位相空間
\newcommand{\Int}[1][]{\mathop{\mathrm{Int}}\nolimits_{#1}}
\newcommand{\Cl}[1][]{\mathop{\mathrm{Cl}}\nolimits_{#1}}
\newcommand{\Bd}[1][]{\mathop{\mathrm{Bd}}\nolimits_{#1}}
\newcommand{\bd}{\mathord{\partial}}
\newcommand{\Fr}[1][]{\mathop{\mathrm{Fr}}\nolimits_{#1}}

\newcommand{\dom}[1][]{\mathop{\mathrm{dom}}\nolimits_{#1}}


\newcommand{\mtan}[1][]{T\hspace{-1.75pt}{#1}}
\newcommand{\mtp}[1][]{{}^{t\!}{#1}} % transpose of the morphism
%\newcommand{\mpb}[1][]{{}^{t\!}{#1}'} % pull-back of the morphism
\newcommand{\mpb}[1][]{{#1}_{\pi}} % pull-back of the morphism
\newcommand{\mdiff}[1][]{{#1}_{d}} % pull-back of the morphism
\newcommand{\mptan}[1][]{{#1}_{\tau}} % pull-back of the morphism

\newcommand{\muinv}[1][]{{#1}_{\mu}^{-1}} % microlocal pull-back of the morphism
\newcommand{\muexinv}[1][]{{#1}_{\mu}^{!}} % microlocal exceptional pull-back of the morphism


%\newcommand{\pb}[1][]{\mathbin{\mathop{\times}\limits_{#1}}}

\newcommand{\cotb}[1][]{T^{\ast}{#1}}
\newcommand{\conm}[2][]{T_{#1}^{\hspace{0.7pt}\ast}{#2}}
\newcommand{\norb}[2][]{T_{#1}{#2}}

\newcommand{\cotz}[1][]{T^{\ast}_{#1}{#1}}

\newcommand{\cotn}[1][]{\dot{T}^{\ast}{#1}}
\newcommand{\conmn}[2][]{\dot{T}_{#1}^{\hspace{0.7pt}\ast}{#2}}


\newcommand{\muhom}{\mu hom}
\newcommand{\muHom}[1][]{\mathrm{Hom}^\mu_{\raise1.5ex\hbox to.1em{}#1}}

\newcommand{\mucat}[2][]{\mathsf{D}^{\mathrm{b}}_{#1}(#2)}
%\newcommand{\mucat}[2][]{\mathsf{D}^{\mathrm{b}}_{#1}(#2)}

\newcommand{\conv}[1][]{\mathbin{\mathop{\circ}\limits_{#1}}}

\newcommand{\torus}{\mathbf{T}}

\newcommand{\diag}[1][]{\Delta_{#1}}
%  \[\mu_{\Delta_{Y}}\mathrm{R}\mathcal{H}om(G,F)\]
\newcommand{\micro}[1]{\mu_{#1}\hspace{0pt}}

\newcommand{\qis}{\stackrel{\text{qis}}{\sim}}%同型右矢印

\newcommand{\Dlcon}{\mathop{\mathsf{D}^{+}_{\rr_{>0}}}\nolimits}
%\newcommand{\Dlcon}{\mathop{\mathsf{D}^{+}_{\rr^+}}\nolimits}
\newcommand{\Dbcon}{\mathop{\mathsf{D}^{\mathrm{b}}_{\rr^+}}\nolimits}

\newcommand{\hfun}[1][]{\mathcal{B}_{#1}} % hyperfunctions
\newcommand{\mfun}[1][]{\mathcal{C}_{#1}} % microfunctions

\newcommand{\wncone}[1][]{C_{#1}} % Whitney normal cones

% ----------------------------
% documenet 
% ----------------------------
% 以下, 本文の執筆スペースです. 
% Your main code must be written between 
% begin document and end document.
% ---------------------------

\begin{document}
\maketitle
\section{余接束における法錐}
次の節では，いろいろな演算のもとでの超局所台の振る舞いを調べる．
結果を定式化するために，
余接束の錐状部分集合に対する新しい演算をいくつか定義しなければならない．
その演算の構成には，第4章第1節で導入した「法錐」の概念と
\(\cotb[X]\)のシンプレクティック構造 (cf.~Appendix) を用いる．
\(X\)を多様体とする．\(\cotb[{\cotb[X]}]\)と\(T\cotb[X]\)をハミルトン同型
\(H\)を用いて\(-H\)により同一視する．
\((x)=(x_1,\dots,x_n)\)を\(X\)の局所座標系とし，
\((x;\xi)\)をそれに対応する\(\cotb[X]\)の座標系とする．



以下の同一視をする．
\[\xymatrix@C=20pt@R=10pt
{
    \cotb[{\norb[M]{X}}]
        \ar@{<->}[r]^-{\varPhi_{\norb[M]{X}}}_-{\sim}
    &
    \cotb[{\conm[M]{X}}]
        \ar@{<->}[r]^-{-H}_-{\sim}
    &
    \norb[{\conm[M]{X}}]{\cotb[X]}\\
    (0,x'',x',0;0,\xi'',\xi',0)
    &
    (0,x'',\xi',0;0,\xi'',-x',0)
    &
    (0,x'',\xi',0;x',0,0,\xi'')
}\]
さらにこれらと\((x',x'';\xi',\xi'')\in \cotb[X]\)を同一視する．

上の座標のもとで\(\cotb[{\conm[M]{X}}]\)には次の2通りの仕方で\(\rrp\)作用が定まる．
\begin{itemize}
    \item 余接空間への作用から定まるもの \((0,x'',\xi',0;0,\xi'',-x',0)\mapsto (0,x'',\xi',0;0,\lambda\xi'',\lambda(-x'),0)\)
    \item 余法束への作用から定まるもの \((0,x'',\xi',0;0,\xi'',-x',0)\mapsto (0,x'',\mu\xi',0;0,\xi'',\mu^{-1}(-x'),0)\)
\end{itemize}

\begin{leftbar}
\begin{LMM}[{\cite[lem.6.2.1]{KS90}}]
    \(X\)を多様体，\(M\)を\(X\)の部分多様体とし，
    \((x)=(x',x'')\)を\(M=\{x'=0\}\)をみたす局所座標とする．
    それに対応する\(\cotb[X]\)の局所座標を\((x',x'';\xi',\xi'')\)で表す．
    \(A\)を\(\cotb[X]\)の錐的集合とする．
    \begin{enumerate}[(i)]
        \item \(\mpb[p]\mdiff[p]^{-1}\wncone[{\conm[M]{X}}](A)=\cotb[M]\cap \wncone[{\conm[M]{X}}](A)\)
        \item \((x''_{\circ};\xi''_{\circ})
            \in \cotb[M]\cap \wncone[{\conm[M]{X}}](A)
        \)となるのは，\(A\)の点列\[
            \left(x'_n,x''_n;\xi'_n,\xi''_n\right)_n
        \]で次の条件\eqref{eq:6.2.4}をみたすものが存在するときである．
        \begin{equation}\label{eq:6.2.4}
            \begin{dcases}
                (x''_n;\xi''_n)\underset{n}{\to}(x''_{\circ};\xi''_{\circ}),\\
                \left\lvert{x'_n}\right\rvert \underset{n}{\to} 0,\\
                \left\lvert{x'_n}\right\rvert \left\lvert{\xi'_n}\right\rvert \underset{n}{\to} 0.
            \end{dcases}
        \end{equation}
        \item \((x''_{\circ};\xi''_{\circ})
            \in \mpb[{\dot{p}}]\mdiff[{\dot{p}}]^{-1}\wncone[{\conmn[M]{X}}](A)
        \)となるのは，\(A\)の点列\[
            \left(x'_n,x''_n;\xi'_n,\xi''_n\right)_n
        \]で\eqref{eq:6.2.4}に加え
        \begin{equation}\label{eq:6.2.5}
            \left\lvert{\xi'_n}\right\rvert \underset{n}{\to} +\infty
        \end{equation}
        をみたすものが存在するときである．
    \end{enumerate}
\end{LMM}
\end{leftbar}

命題 \ref{prp4.1.2}より，\(
    (0,x_0'',\xi_0',0;x_0',0,0,\xi_0'')
    \in \norb[{\conm[M]{X}}]{\cotb[X]}
\)が法錐\(\wncone[{\conm[M]{X}}](A)\)に属するためには，
\(A\times\rrp\)の点列\[
    \left(\left(x_n',x_n'';\xi'_n,\xi''_n\right),c_n\right)_n
\]で\[\begin{dcases}
    \text{\(\conm[M]{X}\)で } (0,x_n'',\xi_n',0)\underset{n}{\to} (0,x_0'',\xi_0',0),\\
    \text{\(\conm[M]{X}\)の法ベクトルとして } c_n(x_n',0,0,\xi_n'')\underset{n}{\to} (x_0',0,0,\xi_0'')
\end{dcases}\]
をみたすものが存在することが必要十分である．

\begin{proof}
    (i) 
    まず，\(\wncone[{\conm[M]{X}}](A)\)が2重錐的 (biconic) であることを示す．
    \(
        (0,x_0'',\xi_0',0;x_0',0,0,\xi_0'')
        \in \norb[{\conm[M]{X}}]{\cotb[X]}
    \)であるとする．
    これを\(\cotb[{\conm[M]{X}}]\)の点
    \(
        (0,x_0'',\xi_0',0;0,\xi_0'',-x_0',0)
        \in \cotb[{\conm[M]{X}}]
    \)とみなす．このとき，
    \((0,x_0'',\xi_0',0;x_0',0,0,\xi_0'')\in\wncone[{\conm[M]{X}}](A)\)
    ならば，
    \[
        (0,x_0'',\xi_0',0;0,\lambda\xi_0'',\lambda(-x_0'),0)
        \leftrightarrow
        (0,x_0'',\xi_0',0;\lambda x_0',0,0,\lambda\xi_0'')
        \in\wncone[{\conm[M]{X}}](A)
    \]かつ
    \[
        (0,x_0'',\mu\xi_0',0;0,\xi_0'',\frac{1}{\mu}(-x_0'),0)
        \leftrightarrow
        (0,x_0'',\mu\xi_0',0;\frac{1}{\mu}x_0',0,0,\xi_0'')
        \in\wncone[{\conm[M]{X}}](A)
    \]であることを示せばよい．

    \((0,x_0'',\xi_0',0;x_0',0,0,\xi_0'')\in\wncone[{\conm[M]{X}}](A)\)
    とする．このとき，\(A\times\rrp\)の点列\(
        \left(
            \left(x_n',x_n'';\xi'_n,\xi''_n\right),
            c_n
        \right)_n
    \)で\begin{align*}
        (0,x_n'',\xi_n',0)
        \underset{n}{\to} (0,x_0'',\xi_0',0),\\
        c_n(x_n',0,0,\xi_n'')
        \underset{n}{\to} (x_0',0,0,\xi_0'')
    \end{align*}
    をみたすものが存在する．
    \(\lambda>0\)を実数とする．
    この点列に対し，\[
        \left(\tilde{x}_n',\tilde{x}_n'';\tilde{\xi}'_n,\tilde{\xi}''_n\right)
        =
        \left(x_n',x_n'';\xi'_n,\xi''_n\right),
        \quad
        \tilde{c}_n=\lambda c_n
    \]
    とおくと，\(\left(
        \left(\tilde{x}_n',\tilde{x}_n'';\tilde{\xi}'_n,\tilde{\xi}''_n\right),
        \tilde{c}_n
    \right)_n\)は\(A\times \rrp\)の点列であり，
    \begin{align*}
        (0,\tilde{x}_n'',\tilde{\xi}_n',0)
        &=(0,x_n'',\xi_n',0)\\
        &\to(0,x_0'',\xi_0',0),\\
        \tilde{c}_n(\tilde{x}_n',0,0,\tilde{\xi}_n'')
        &=(\lambda c_n x_n',0,0,\lambda c_n \xi_n'')\\
        &\to(\lambda x_0',0,0,\lambda\xi_0'')
    \end{align*}
    である．したがって，
    \(
        (0,x_0'',\xi_0',0;\lambda x_0',0,0,\lambda\xi_0'')
        \in\wncone[{\conm[M]{X}}](A)
    \)である．

    次に，\(\mu>0\)を実数とし，\[
        \left(\check{x}_n',\check{x}_n'';\check{\xi}'_n,\check{\xi}''_n\right)
        =
        \left(x_n',x_n'';\mu\xi'_n,\mu\xi''_n\right),
        \quad 
        \check{c}_n=\frac{1}{\mu} c_n
    \]
    とおくと，\(A\)が錐状であることから\(\left(
        \left(\check{x}_n',\check{x}_n'';\check{\xi}'_n,\check{\xi}''_n\right),
        \check{c}_n
    \right)_n\)は\(A\times \rrp\)の点列であり，
    \begin{align*}
        (0,\check{x}_n'',\check{\xi}_n',0)
        &=(0,x_n'',\mu\xi_n',0)\\
        &\to(0,x_0'',\mu\xi_0',0),\\
        \check{c}_n(\check{x}_n',0,0,\check{\xi}_n'')
        &=\left(\frac{1}{\mu}c_n x_n',0,0,\frac{1}{\mu}c_n\cdot \mu\xi_n''\right)\\
        &=\left(\frac{1}{\mu}c_n x_n',0,0,c_n\xi_n''\right)\\&\to\left(\frac{1}{\mu} x_0',0,0,\xi_0''\right)
    \end{align*}
    である．よって\(
        (0,x_0'',\mu\xi_0',0;\mu^{-1}x_0',0,0,\xi_0'')
        \in\wncone[{\conm[M]{X}}](A)
    \)である．
\end{proof}

\appendix
\section*{本編で用いる命題}
\begin{leftbar}
\begin{PRP}[{\cite[Prop. 4.1.2]{KS90}}]\label{prp4.1.2}
    (ii) 
    点列
\end{PRP}
\end{leftbar}
%\section{aa}
%===============================================
% 参考文献スペース
%===============================================
\begin{thebibliography}{20} 
  \bibitem[dR]{dR} 
  de Rham, \textit{Vari\'et\'es diff\'erentiables}, Hermann, 1953.
  和訳：ド・ラーム, 微分多様体, 東京図書, 1974.
  \bibitem[GKS12]{GKS12} St\'ephane Guillermou, 
  Masaki Kashiwara, Pierre Schapira, 
  \textit{Sheaf Quantization of Hamiltonian Isotopies and Applications to Nondisplaceability Problems}, 
  Duke. Math. J. 161, No. 2, pp.201--245, 2012.  
  \bibitem[GP74]{GP74} Victor Guillemin, Alan Pollack, 
    \textit{Differential Topology}, 
    Prentice-Hall, 1974.
  \bibitem[KS90]{KS90} Masaki Kashiwara, Pierre Schapira, 
    \textit{Sheaves on Manifolds}, 
    Grundlehren der Mathematischen Wissenschaften, 292, Springer, 1990.
  \bibitem[Hat89]{Hat89} 服部晶夫, 多様体 増補版, 岩波全書 288, 岩波書店, 1989.
  %\bibitem[Og02]{Og02} 小木曽啓示, 代数曲線論, 朝倉書店, 2022.
  \bibitem[Hor63]{Hor63} Lars H\"ormander, 
  \textit{Linear Pertial Differential Operators}, Fourth Printing, 
  Grundlehren der Mathematischen Wissenschaften, 116, Springer, 1963.
  \bibitem[Hor71]{Hor71} Lars H\"ormander, 
  \textit{Fourier Integral Operators I}, 
  Acta Math. 127, 79--183, (1971).
  \bibitem[Hor89]{Hor89} Lars H\"ormander, 
  \textit{The Analysis of Linear Pertial Differential Operators I}, Second Edition,
  Grundlehren der Mathematischen Wissenschaften, 256, Springer, 1989.
  \bibitem[Hor85]{Hor85} Lars H\"ormander, 
  \textit{The Analysis of Linear Pertial Differential Operators III}, 
  Grundlehren der Mathematischen Wissenschaften, 274, Springer, 1985.
  \bibitem[今野13]{今野13} 今野宏, 微分幾何学, 東京大学出版会, 2013.
  \bibitem[Roc]{Roc} Rockaffeller, \textit{Convex Analysis},. 
  \bibitem[Sch85]{Sch85} Pierre Schapira, 
  \textit{Microdifferential Systems in the Complex Domain}, 
  Grundlehren der Mathematischen Wissenschaften, 269, 
  Springer, 1985.
  \bibitem[Sch88]{Sch88} Pierre Schapira, 
  \textit{Microfunctions for Boundary Value Problems}, 
  in Algebraic Analysis, Volume II, 1988.
  \bibitem[ST92]{ST92} Pierre Schapira, Nobuyuki Tose, 
  \textit{Morse Inequalities for \(\rr\)-constructible Sheaves}, 
  Adv. Math. 93, pp.1--8, 1992. 
  \bibitem[Zwo12]{Zwo12} Maciej Zworski,
  \textit{Semiclsaaical Analysis}, 
  Graduate Studies in Mathematics, 138, American Mathematical Society, 2012. 
\end{thebibliography}

%===============================================




\end{document}
